% Originally: content/week-02.tex, \section{Continuity} (Corsin Week 2, Friday)

% Source: Corsin Nick/Class Notes/Week 2.pdf, pp. 1--13
\section{Continuity}
\label{sec:continuity}

\begin{definition}[Continuous map]
\label{def:continuous}
For $(X, d_X)$ and $(Y, d_Y)$ metric spaces, a function $f : X \to Y$ is called \newterm{continuous}
\germanfor{stetig} if one of the following equivalent conditions holds:
\begin{enumerate}[label=\textbf{(\alph*)}]
  \item \textbf{($\varepsilon$--$\delta$):} $\forall x_0 \in X, \varepsilon > 0 \ \exists \delta > 0$ such that for all $x \in X$:
    \[d(x_0, x) < \delta \implies d(f(x_0), f(x)) < \varepsilon\]
  \item \textbf{(Sequentially continuous):} for any sequence $(x_n)_{n=0}^{\infty}$ in $X$ with
    limit $x \in X$, $\lim_{n\to\infty} f(x_n) = f(x)$ (and the limit exists).
  \item \textbf{(Topological):} for all $U \subseteq Y$ open, $f^{-1}(U) \subseteq X$ is open
    (\cref{def:open_closed}).
\end{enumerate}
\end{definition}

\begin{ainote}
Corsin writes $d(x_0,x) < \varepsilon$ in the ($\varepsilon$--$\delta$) hypothesis above, which
must read $\delta$, since otherwise the quantified $\delta$ is never used at all. This is a slip
of the pen, and it is corrected in \cref{def:continuous}.
\end{ainote}

% Generator: Claude Opus 5 (High) --- \cref{def:continuous} asserts the equivalence of three
% conditions without proving it; the cycle is short and uses only results already available.
\begin{proposition}[The three definitions of continuity agree]
\label{prop:continuity_equivalences}
The conditions \textbf{(a)}, \textbf{(b)} and \textbf{(c)} of \cref{def:continuous} are
equivalent, so \cref{def:continuous} is well posed.
\end{proposition}

\begin{proof}
We prove the cycle
$\text{\textbf{(a)}} \Rightarrow \text{\textbf{(b)}} \Rightarrow \text{\textbf{(c)}}
\Rightarrow \text{\textbf{(a)}}$.

\textbf{(a) $\Rightarrow$ (b).} Let $x_n \to x$ and let $\varepsilon>0$. Pick $\delta>0$ as in
\textbf{(a)} for the point $x$, then $N$ with $d_X(x_n,x)<\delta$ for $n \geq N$. For those $n$,
$d_Y(f(x_n),f(x)) < \varepsilon$, so $f(x_n) \to f(x)$.

\textbf{(b) $\Rightarrow$ (c).} Let $U \subseteq Y$ be open and suppose $f^{-1}(U)$ were not.
By the shrinking-ball construction (\cref{rem:shrinking_ball_trick}) there are
$x \in f^{-1}(U)$ and points $x_n \in B_{1/n}(x)$ with $x_n \notin f^{-1}(U)$. Then $x_n \to x$,
so $f(x_n) \to f(x) \in U$ by \textbf{(b)}, and since $U$ is open,
\cref{prop:open_closed_via_sequences}\textbf{(a)} forces $f(x_n) \in U$ for large $n$. That is,
$x_n \in f^{-1}(U)$, which is a contradiction.

\textbf{(c) $\Rightarrow$ (a).} Fix $x_0 \in X$ and $\varepsilon>0$. The ball
$B_\varepsilon(f(x_0)) \subseteq Y$ is open (\cref{ex:ai_open_balls_are_open}), so
$f^{-1}\big(B_\varepsilon(f(x_0))\big)$ is open by \textbf{(c)} and contains $x_0$. Hence there
is $\delta>0$ with $B_\delta(x_0) \subseteq f^{-1}\big(B_\varepsilon(f(x_0))\big)$, which says
precisely that $d_X(x_0,x)<\delta$ implies $d_Y(f(x_0),f(x))<\varepsilon$.
\end{proof}

% Generator: Claude Opus 5 (High) --- FIG-05-01. The chapter defines continuity three ways and
% had no picture at all; its only figure was the Lipschitz cone, two sections later.
%
% GEOMETRY. Nothing here is a computation, so there is no arithmetic to check; what matters is
% that the drawing asserts only true things. Two deliberate choices:
%   * The image of B_delta(x_0) is drawn as an irregular blob strictly INSIDE B_epsilon(f(x_0)),
%     not as a disc and not filling it. Continuity says the image is contained in the target
%     ball; it says nothing about the image being round, being centred, or exhausting the ball.
%     A figure drawing f(B_delta) as a concentric disc would assert all three.
%   * The blob touches neither the boundary of B_epsilon nor f(x_0)'s dot, since the inclusion
%     is into the OPEN ball and f(x_0) is interior to the image, not on its edge.
% The same picture reads condition (c) backwards: B_delta sits inside the preimage of the open
% ball B_epsilon, which is the containment the proof of (c) => (a) produces.
\begin{center}
\begin{tikzpicture}[>=Latex, scale=1.0]
  % ---- the source space ----
  \draw[gray!55] plot[smooth cycle, tension=0.75] coordinates
    {(-2.5,1.15) (-0.85,1.5) (0.15,0.55) (-0.35,-1.2) (-1.9,-1.45) (-2.95,-0.25)};
  \node[gray!55!black, font=\small] at (-2.72,1.28) {$X$};
  \fill[MidnightBlue!8] (-1.45,0.05) circle (0.62);
  \draw[MidnightBlue!70, thick] (-1.45,0.05) circle (0.62);
  \fill[MidnightBlue!70] (-1.45,0.05) circle (1.4pt);
  \node[MidnightBlue!70, font=\scriptsize, anchor=north east] at (-1.5,-0.02) {$x_0$};
  \node[MidnightBlue!70, font=\scriptsize, anchor=south] at (-1.45,0.70) {$B_\delta(x_0)$};

  % ---- the map ----
  \draw[->, thick] (0.55,0.05) -- (1.95,0.05) node[midway, above, font=\small] {$f$};

  % ---- the target space ----
  \draw[gray!55] plot[smooth cycle, tension=0.75] coordinates
    {(2.6,1.3) (4.4,1.5) (5.5,0.4) (4.9,-1.3) (3.3,-1.5) (2.35,-0.3)};
  \node[gray!55!black, font=\small] at (5.35,1.15) {$Y$};
  \fill[BrickRed!7] (3.95,0.05) circle (0.78);
  \draw[BrickRed!80, thick] (3.95,0.05) circle (0.78);
  \node[BrickRed!80, font=\scriptsize, anchor=south] at (3.95,0.86)
    {$B_\varepsilon(f(x_0))$};

  % The image: irregular, strictly inside, not centred. All three are deliberate.
  \fill[OliveGreen!14] plot[smooth cycle, tension=0.8] coordinates
    {(3.72,0.42) (4.25,0.3) (4.38,-0.12) (3.95,-0.42) (3.55,-0.18)};
  \draw[OliveGreen, thick] plot[smooth cycle, tension=0.8] coordinates
    {(3.72,0.42) (4.25,0.3) (4.38,-0.12) (3.95,-0.42) (3.55,-0.18)};
  \fill[BrickRed!80] (3.95,0.02) circle (1.4pt);
  % Anchored below the dot and inside the blob. Anchored east it extended leftwards across the
  % blob's own outline at x = 3.55, so the label sat half in and half out of the region it names.
  \node[BrickRed!80, font=\scriptsize, anchor=north] at (3.95,-0.04) {$f(x_0)$};
  \node[OliveGreen, font=\scriptsize, anchor=north west] at (4.35,-0.35)
    {$f\big(B_\delta(x_0)\big)$};
\end{tikzpicture}
\end{center}

Condition \textbf{(a)} is the statement that the green region fits inside the red ball, and that
a $\delta$ making it fit can be found however small the red ball is made. Condition \textbf{(c)}
is the same picture read from right to left: $B_\delta(x_0)$ lies inside
$f^{-1}\big(B_\varepsilon(f(x_0))\big)$, which is exactly the containment produced in the
\textbf{(c)} $\Rightarrow$ \textbf{(a)} step above. Note what the drawing does not claim: the
image need not be round, need not be centred at $f(x_0)$, and need not fill the ball.

\begin{remark}[Which of the three to reach for]
Condition \textbf{(c)} is the only one that never mentions the metric. That is why it, and not the
$\varepsilon$--$\delta$ formulation, is what survives into general topology
(\cref{sec:structured_spaces}). Condition \textbf{(b)} is usually the fastest way to prove a
function is \emph{dis}continuous, since a single badly chosen sequence suffices; this is exactly
how \cref{ex:partial_derivatives_not_continuous} works. Condition \textbf{(a)} is the
one to use when a quantitative modulus is wanted, as in
\cref{def:uniformly_continuous} immediately below.
\end{remark}

% Sonnet 5 (Medium)
\begin{definition}[Uniform continuity]
\label{def:uniformly_continuous}
For $(X,d_X)$ and $(Y,d_Y)$ metric spaces, a function $f : X \to Y$ is \newterm{uniformly
continuous} \germanfor{gleichm\"assig stetig} if
\[\forall \varepsilon > 0 \ \exists \delta > 0 \ \forall x, x' \in X : \quad d_X(x,x') < \delta \ \text{ implies } \ d_Y(f(x),f(x')) < \varepsilon.\]
\end{definition}

\begin{importantremark}[Where the quantifiers sit]
\label{rem:continuity_vs_uniform_continuity}
Compare \cref{def:continuous} \textbf{(a)} with \cref{def:uniformly_continuous}: the
\emph{formulas} differ only in the order of two quantifiers, and that order is the whole
content.
\begin{itemize}
  \item \textbf{Continuity:} $\forall x_0\ \forall\varepsilon\ \exists\delta$. The point $x_0$
    is chosen \emph{before} $\delta$, so $\delta$ may depend on it, and a different $\delta$ is
    allowed at every point.
  \item \textbf{Uniform continuity:} $\forall\varepsilon\ \exists\delta\ \forall x$. Now
    $\delta$ is chosen \emph{before} seeing the point, so one and the same $\delta$ must work
    everywhere at once.
\end{itemize}
Uniform continuity is therefore strictly stronger. The standard witness is
$f(x) := 1/x$ on $(0,1)$: it is continuous, but near $0$ the graph steepens without bound, so
the $\delta$ that works at $x_0 = \tfrac12$ is hopeless at $x_0 = \tfrac{1}{1000}$, and no
single $\delta$ survives. \Cref{prop:compact_implies_uniformly_continuous} below says that on a
\emph{compact} space this can never happen. The gap between the two notions closes exactly when
there is no room left to run off towards a bad point.
\end{importantremark}

\begin{example}[$f(x) = x^2$]
Let $f : \mathbb{R} \to \mathbb{R}, x \mapsto x^2$. Let $a < b \in \mathbb{R}$. Then
\begin{itemize}
  \item for $a < 0 < b$: $\quad f^{-1}((a,b)) = (-\sqrt{b}, \sqrt{b})$, which is open.
  \item for $a < b < 0$: $\quad f^{-1}((a,b)) = \emptyset$.
  \item for $0 < a < b$: $\quad f^{-1}((a,b)) = (-\sqrt{b}, -\sqrt{a}) \cup (\sqrt{a}, \sqrt{b})$.
\end{itemize}
\textit{Why is it enough to check only these cases?} Hint:
$f^{-1}(U_1 \cup U_2) = f^{-1}(U_1) \cup f^{-1}(U_2)$.
\end{example}

% Source: Jérôme Paschoud/Notizen/Woche 4.1 Normen, das Differential.pdf, p. 2
% Generator: Gemini 3.1 Pro (High)
\begin{example}[Limits in $\mathbb{R}^2$ depend on the path]
\label{ex:limits_in_r2}
When computing limits in $\mathbb{R}^n$, one must ensure that the limit exists and is identical along \emph{all} possible paths approaching the point. For example, consider the function $f: \mathbb{R}^2 \setminus \{(0,0)\} \to \mathbb{R}$ given by $f(x,y) = \frac{x^2}{x^2+y^2}$.
If we approach the origin along the $x$-axis (setting $y=0$ first), we get
\[
\lim_{x \to 0} f(x, 0) = \lim_{x \to 0} \frac{x^2}{x^2} = 1.
\]
However, if we approach along the $y$-axis (setting $x=0$ first), we get
\[
\lim_{y \to 0} f(0, y) = \lim_{y \to 0} \frac{0}{y^2} = 0.
\]
Since these limits disagree, $\lim_{(x,y) \to (0,0)} f(x,y)$ does not exist. A robust method to check limits at the origin uniformly across all directions is to switch to \newterm{polar coordinates} ($x = r\cos\varphi, y = r\sin\varphi$). If the limit as $r \to 0$ is independent of $\varphi$, then the function converges. (This is explored further in \cref{ex:3.2}).
\end{example}

\begin{question}
Let $(X, d_X)$ and $(Y, d_Y)$ be metric spaces, where $d_Y$ is the discrete metric. Is
\emph{every} function $f : (X,d_X) \to (Y,d_Y)$ continuous?
\end{question}

\begin{answer}
\textbf{No.} For example, take the spaces $X = Y = \mathbb{R}$, $d_X$ the standard metric, and
$f = \id : X \to Y, z \mapsto z$. Then $\{0\} \subseteq Y$ is open, but
$\{0\} = \id^{-1}(\{0\})$ is not open in $(X, d_X)$.
\end{answer}

\begin{question}
And what if the roles are reversed, with $d_X$ the discrete metric and $d_Y$ arbitrary?
\end{question}

\begin{answer}
\textbf{Any} subset of $X$ is open; therefore, $f^{-1}(U)$ is open for any open $U \subseteq Y$. So
every such $f$ is continuous.
\end{answer}

% The old chapter continued from here into "Completeness"; that topic is now
% \cref{ch:completeness}, whose opening paragraph carries this transition.


% --- exercises relocated here from the dissolved sheets ---

% Originally: content/exercise-sheets/week-02.tex, problem 2.6
% Source: exercises/Ex2_Analysis2_eng.pdf, p. 1

\begin{exercise}[Continuity of the distance]
\label{ex:2.6}
Let $(X,d)$ be a metric space, and endow $X \times X$ with the product distance
$d_2(x,y) := d(x_1,y_1) + d(x_2,y_2)$. Show that the distance function
$d : X \times X \to \mathbb{R}$ is continuous with respect to $d_2$ and, more precisely, that it
is 1-Lipschitz. Is $d$ also continuous with respect to the distance
$d_1(x,y) := \max\{d(x_1,y_1), d(x_2,y_2)\}$? Is $d$ also Lipschitz continuous with respect to the
distance $d_3(x,y) := \sqrt{d(x_1,y_1)^2 + d(x_2,y_2)^2}$? (Notation consistent with
\cref{ex:2.5}.)
\exinfo{This exercise is Problem 2.6 of Problem Sheet 2, Analysis II, Spring Semester 2026.}
\end{exercise}

% Originally: content/exercise-sheets/week-02.tex, problem 2.7
% Source: exercises/Ex2_Analysis2_eng.pdf, p. 1

\begin{exercise}[Continuity of the composition]
\label{ex:2.7}
Let $X, Y, Z$ be metric spaces and let $f : X \to Y$, $g : Y \to Z$ be continuous functions. Show
that $g \circ f : X \to Z$ is continuous using at least two of the three equivalent definitions of
continuity seen in class.
\exinfo{This exercise is Problem 2.7 of Problem Sheet 2, Analysis II, Spring Semester 2026.}
\end{exercise}

\begin{remark}
The three equivalent definitions of continuity referenced in \cref{ex:2.7} are exactly the ones
stated in \cref{sec:continuity} below (\(\varepsilon\)--\(\delta\), sequential, topological).
\end{remark}

% Originally: content/exercise-sheets/week-03.tex, problem 3.2
% Source: exercises/Ex3_Analysis2_eng.pdf

\begin{exercise}[Problems at the origin]
\label{ex:3.2}
Show that there is no continuous function $f : \mathbb{R}^2 \to \mathbb{R}$ such that
$f(x,y) = xy/(x^2+y^2)$ for all $(x,y) \neq (0,0)$. On the other hand, show that there is exactly
one continuous function $g : \mathbb{R}^2 \to \mathbb{R}$ such that
$g(x,y) = xy/\sqrt{x^2+y^2}$ for all $(x,y) \neq (0,0)$.

\exhint[Official hint]{If a function is continuous at $0$ and $x_k \to 0$ and $y_k \to 0$, then
$f(x_k)$ and $f(y_k)$ must converge to the same value.}
\exinfo{This exercise is Problem 3.2 of Problem Sheet 3, Analysis II, Spring Semester 2026.}
\end{exercise}

% Quelle: Jérôme Paschoud/Notizen/Woche 2.2 Topologie und Stetigkeit.pdf, p. 2
\begin{exercise}[Distance to a point is continuous]
\label{ex:distance_to_point_continuous}
Let $(X, d)$ be a metric space and fix $a \in X$. Show that the function $f : X \to [0,\infty)$ defined by $f(x) := d(x, a)$ is continuous.
\end{exercise}

% Generator: Gemini 3.6 Pro
\begin{exercisesolution}[Proof of \cref{ex:distance_to_point_continuous}]
We use $\varepsilon$--$\delta$ continuity. Let $x \in X$ and $\varepsilon > 0$. We need to find $\delta > 0$ such that for all $x' \in X$ with $d(x, x') < \delta$, we have $|f(x) - f(x')| < \varepsilon$.
Using the reverse triangle inequality, we have
\[|f(x) - f(x')| = |d(x, a) - d(x', a)| \leq d(x, x').\]
Therefore, we can simply choose $\delta := \varepsilon$. Then $d(x, x') < \delta$ directly implies $|f(x) - f(x')| < \varepsilon$, meaning $f$ is continuous at $x$. Since $x$ was arbitrary, $f$ is continuous everywhere.
\end{exercisesolution}

% Supplement: Analysis_II_Script_v1.pdf, p. 19
\begin{corollary}[Closed sets are exactly the zero sets of a distance function]
\label{cor:closed_iff_zero_of_distance}
Let $(X,d)$ be a metric space and $\emptyset \neq E \subseteq X$. The function
$f_E(x) := d(x,E) = \inf_{z\in E}d(x,z)$ is $1$-Lipschitz, which is \cref{ex:2.3}\textbf{(c)},
proved there by the same reverse-triangle-inequality argument as
\cref{ex:distance_to_point_continuous} above. Moreover,
\[E \text{ is closed} \quad\Longleftrightarrow\quad E = \{x \in X : f_E(x) = 0\}.\]
\end{corollary}

\begin{proof}
The inclusion $E \subseteq \{f_E = 0\}$ always holds, since $z \in E$ gives
$f_E(z) \leq d(z,z) = 0$. So only the reverse inclusion carries content, and it is exactly where
closedness is needed.

\textbf{($\Rightarrow$)} Let $E$ be closed and $f_E(x) = 0$. By definition of the infimum, choose
$z_n \in E$ with $d(x,z_n) \to 0$; then $z_n \to x$, and closedness of $E$ gives $x \in E$
(\cref{prop:open_closed_via_sequences}\textbf{(b)}).

\textbf{($\Leftarrow$)} Suppose $E = \{f_E = 0\}$ and let $(z_n)$ be a sequence in $E$ with
$z_n \to x$. Then $f_E(x) \leq d(x,z_n) \to 0$, so $f_E(x) = 0$, i.e.\ $x \in E$ by hypothesis.
\Cref{prop:open_closed_via_sequences}\textbf{(b)} then gives that $E$ is closed.
\end{proof}

\begin{remark}
This sharpens \cref{ex:2.3}\textbf{(a)}, which showed only one direction ($E$ closed
$\Rightarrow$ $f_E > 0$ off $E$) for a general subset. It is exactly closedness, and nothing
weaker, that lets $f_E$ detect membership in $E$ by vanishing.
\end{remark}

% Quelle: Jérôme Paschoud/Notizen/Woche 2.2 Topologie und Stetigkeit.pdf, p. 3
\begin{exercise}[Evaluation map is continuous]
\label{ex:evaluation_map_continuous}
Let $X := C^0([0,1])$ equipped with the supremum metric $d(f,g) := \sup_{t \in [0,1]} |f(t) - g(t)|$. Define the function $F : X \to \mathbb{R}$ by $F(f) := f(0)$. Show that $F$ is continuous.
\end{exercise}

% Generator: Gemini 3.6 Pro
\begin{exercisesolution}[Proof of \cref{ex:evaluation_map_continuous}]
Let $f \in X$ and $\varepsilon > 0$. Choose $\delta := \varepsilon$. If $g \in X$ satisfies $d(f, g) < \delta$, then by the definition of the supremum metric, we have $|f(t) - g(t)| < \delta$ for all $t \in [0,1]$.
In particular, for $t = 0$, we get
\[|F(f) - F(g)| = |f(0) - g(0)| \leq \sup_{t \in [0,1]} |f(t) - g(t)| = d(f, g) < \delta = \varepsilon.\]
Thus $F$ is continuous at $f$. Since $f$ was arbitrary, $F$ is continuous on all of $X$.
\end{exercisesolution}
